\originalsection{作业 10: 惯量与力矩}
\begin{exercise}
原题 1. (a) 均匀薄圆盘 $M,R$, 求垂直盘面并分别过质心、轮缘的惯量. (b) 钢垫圈内半径 13.5 mm、外半径 31.0 mm、厚 4.0 mm、密度 $7.8\times10^3$ kg/m$^3$, 求垂直盘面质心轴惯量.
\end{exercise}
\begin{solution}
(a) 第\ref{ch:rotation}章积分给 $I_C=MR^2/2$, 平行轴得 $I_{\rm rim}=3MR^2/2$. (b) $M=\rho\pi d(r_o^2-r_i^2)$, $I=M(r_o^2+r_i^2)/2=4.36\times10^{-5}$ kg m$^2$. 所有毫米先换成米, 不能把厚度省略出质量公式.
\end{solution}
\begin{exercise}
原题 2. 固定复合轮由同轴刚连的均匀盘组成: $M_1,R$ 与 $M_2,R/2$. 绳绕小盘再过无摩擦柱, 末端挂 $M_b$. 轴无摩擦, 从静止释放. 求下落距离 $d$ 所需时间.
\end{exercise}
\begin{solution}
轮惯量 $I=M_1R^2/2+M_2R^2/8$. 绕绳半径 $r=R/2$, 方程 $M_bg-T=M_ba$, $Tr=I a/r$. 故 $a=M_bg/(M_b+2M_1+M_2/2)$, $t=\sqrt{2d/a}$. 无摩擦柱仅改变绳方向, 不改变张力.
\end{solution}
\begin{exercise}
原题 3. 均匀水平杆长 2.0 m、质量 4.0 kg, 左端铰墙, 右端用与杆夹 $\beta=30^\circ$ 的无质量缆绳上拉. (a) 缆绳张力. (b) 墙对杆支点力相对杆的角度. (c) 支点力大小.
\end{exercise}
\begin{solution}
力矩平衡给 $T=mg/(2\sin\beta)=39.24$ N. 支点力分量为 $H_x=T\cos\beta$, $H_y=mg-T\sin\beta=mg/2$, 所以 (b) 角为 $30^\circ$, (c) 大小 39.24 N, 方向右上. 杆长在这个静力关系中约去.
\end{solution}
\begin{exercise}
原题 4. 人质量 $m$ 站倾角 $\alpha$ 坡上, 两足沿坡相距 $d$. 质心在两足中点法向上方 $h$. 静摩擦足够不滑. (a) 求两脚法向力, $N_1>N_2$. (b) 何脚距使上坡脚法向力刚为零、开始翻倒?
\end{exercise}
\begin{solution}
法向合力 $N_d+N_u=mg\cos\alpha$. 关于质心力矩平衡给 $(d/2)(N_d-N_u)=hmg\sin\alpha$. 故下坡脚 $N_1=mg(\cos\alpha/2+h\sin\alpha/d)$, 上坡脚 $N_2=mg(\cos\alpha/2-h\sin\alpha/d)$. (b) $d=2h\tan\alpha$, 维持双脚接触需 $d$ 不小于此值. 摩擦各脚的分配可不唯一, 但总切向平衡固定上述力矩关系.
\end{solution}
\begin{exercise}
原题 5. 质量 $m$, 半径 $R$ 的圆柱在对称 V 槽内原地逆时针转动, 角速度保持常数. 两侧各与水平夹 $45^\circ$, 摩擦系数 $\mu$. 求需施加的外力矩. 原题标题称 rolling, 图实际为两接触处滑动的原地转动.
\end{exercise}
\begin{solution}
设两侧法向力 $N_L,N_R$, 摩擦分别沿各接触速度反向. 水平平衡给 $N_L-N_R=\mu(N_L+N_R)$, 竖直平衡给 $N_L+N_R+\mu(N_L-N_R)=\sqrt2mg$. 故总法向力 $S=\sqrt2mg/(1+\mu^2)$, 两摩擦力矩均顺时针. 外力矩须逆时针, 大小 $\tau=\sqrt2\mu mgR/(1+\mu^2)$. 保持两面均接触还需 $0\le\mu\le1$; 更大系数下这个无其他支承的双接触匀转模型不可实现.
\end{solution}
\begin{exercise}
原题 6. 质量 $m_p$、半径 $R$ 的固定滑轮惯量 $m_pR^2/2$, 无摩擦轴. 轻绳绕轮连倾角 $\beta$ 的光滑斜面块 $m$. 从静止释放, 求块下滑 $d$ 所需时间.
\end{exercise}
\begin{solution}
$mg\sin\beta-T=ma$, $TR=(m_pR^2/2)a/R$. 得 $a=mg\sin\beta/(m+m_p/2)$, $t=\sqrt{2d(m+m_p/2)/(mg\sin\beta)}$.
\end{solution}
